% conclusion.tex
\section{Conclusion}
\label{sec:conclusion}

We took the common argument that benchmark contamination is a privacy warning sign, restated it as
four explicit assumptions, and tested the first of them on models whose training corpus is public.
The question was whether the contamination signal an auditor can cheaply compute predicts concrete
extraction once the model's per-item loss is held fixed. Using a pre-registered partial-correlation
and mediation analysis that controls for raw per-item loss, we found that it does, but only through
loss: the calibrated reference-free detectors (Min-K\%, Min-K\%++, zlib) add no positive independent
predictive value beyond loss. The two negative residuals we observe are consistent with statistical
suppression under near-collinearity rather than with inverse prediction, and we claim only the
absence of positive value. The result is robust to a non-linear loss control and to deduplication, and is not a frequency
or zero-inflation artifact. The practical message is a measurement-validity bound: the detectors
optimized for membership inference carry little information about \emph{which} items leak that loss
does not already carry, so an auditor should measure loss and extractability directly. We claim no
new detector or metric. The contributions are the explicit assumption chain of
Section~\ref{sec:chain} and the controlled, pre-registered measurement of incremental value. These findings are preliminary, on
the smallest Pythia model. The immediate next step, and the design target of our released
pipeline, is the GPU-scale replication across model sizes, where memorization and extraction are
expected to strengthen, and where assumptions A3 and A4 of Section~\ref{sec:chain} could be tested
rather than assumed, and where the question of whether calibrated detectors gain
independent leakage-predictive value at scale can be settled.
